Composite visualizations combine multiple charts through operations such as
concatenation, layering, nesting, and faceting to reveal relationships across
views. Existing authoring tools often present a trade-off between flexibility
and authoring effort: chart-level tools simplify construction but are
constrained by predefined chart and composition structures, whereas
element-level tools offer finer control at the cost of low-level graphical
operations and encoding specifications. We present {\system}, an interactive
authoring tool that treats charts as reusable and editable blocks and uses
explicit visual structure as a reference for authoring. To ground this
abstraction in users' mental models, we conducted a formative study with
visualization experts to understand how composite visualizations are
decomposed into meaningful chart-level units, informing the granularity and
family organization of blocks in {\system}. Authors construct compositions by
spatially arranging and directly combining blocks, with layout serving as a
reference for explicit and editable composition relationships. Each block also
provides a visible visual form with bounded dimensional capacity; when an
additional dimension exceeds this capacity, {\system} exposes related block
forms or composition-level adaptation while preserving compatible design
choices. Through case studies, a comparative user study of implementing a
predefined composite visualization, and an exploratory study spanning broader
chart types, coordinate systems, and composition structures, we show that
{\system} supports composite visualization authoring through a consistent,
learnable, and predictable workflow.
